\section{Introduction}
\label{sec:introduction}

Human memory is not a passive recording device; it is a dynamic, adaptive, and constructive process. We remember what matters, forget what is trivial, and consolidate fragmented experiences into coherent knowledge during sleep. In stark contrast, current memory-augmented Large Language Models (LLMs) implicitly adopt a \textbf{``FIFO or Accumulation-Only''} philosophy. Systems like LangChain~\cite{langchain}, MemGPT~\cite{memgpt}, and Generative Agents~\cite{generative_agents} treat memory as a static store, indiscriminately accumulating context until capacity limits are reached. This leads to a \textbf{``Semantic Quicksand''} effect: as the memory store grows, retrieval precision paradoxically deteriorates because the system lacks a \textit{cognitive filter} to distinguish transient noise from persistent knowledge. Consequently, these systems suffer from three critical failures: (1) \textbf{Memory Bloat}, where token costs scale linearly with interaction length; (2) \textbf{Semantic Stagnation}, where specific events fail to evolve into generalized knowledge; and (3) \textbf{Poor Personalization}, where a ``one-size-fits-all'' retention policy ignores individual user differences.

While recent works have attempted to address these issues, significant gaps remain. Generative Agents~\cite{generative_agents} introduce reflection mechanisms, but they are computationally prohibitive for online dialogue and lack adaptive parameters. Machine Unlearning~\cite{sekhar2021unlearning} focuses on erasing data from model weights, which is too mechanical and lacks the nuance of psychological forgetting. There is a lack of a principled cognitive framework for \textbf{online, personalized, and evolutional} memory management that bridges the gap between human cognition and LLM architecture.

To bridge this gap, we introduce \textbf{Mem0-Cognitive}, a novel framework that reframes memory management not as a storage problem, but as an \textbf{Active Inference} process. Inspired by cognitive psychology, our system operationalizes three core principles: \textbf{Emotion-Weighted Forgetting}, \textbf{Offline Sleep Consolidation}, and \textbf{Meta-Cognitive Adaptation}. Unlike static baselines, Mem0-Cognitive continuously predicts the future utility of memory traces, dynamically pruning redundancy while reinforcing salient information.

\textbf{Our Contributions} are fourfold:
\begin{enumerate}
    \item \textbf{A Cognitively-Inspired Framework}: We introduce the Mem0-Cognitive framework, which integrates emotion-weighted Ebbinghaus decay, offline schema induction via sleep consolidation, and a meta-cognitive learner for personalized parameter adaptation.
    \item \textbf{Algorithmic Formalization}: We derive a novel \textit{Affective Retention Score} that incorporates a \textbf{Salience Gate} mechanism, enabling dynamic memory pruning without catastrophic forgetting of emotionally charged events.
    \item \textbf{Benchmark \& Evaluation}: We release \textbf{CognitiveBench}, a long-context dialogue simulator designed to probe memory retention vs. noise trade-offs over 1000+ turns, alongside evaluations on the LoCoMo benchmark.
    \item \textbf{Empirical Findings}: Extensive experiments demonstrate a \textbf{55\% reduction in token overhead} and a \textbf{29\% increase in critical fact retention} compared to state-of-the-art baselines, with sub-50ms inference latency overhead. Ablation studies confirm the distinct contributions of each cognitive module.
\end{enumerate}
